\documentclass[12pt]{article}
\usepackage{amsmath,amssymb,amsthm}
\usepackage{aliascnt}
\usepackage{hyperref}
\newtheorem{theorem}{Theorem}[section]
\newaliascnt{lemma}{theorem}
\newtheorem{lemma}[lemma]{Lemma}
\aliascntresetthe{lemma}
\providecommand{\lemmaautorefname}{Lemma}
\theoremstyle{definition}
\newaliascnt{definition}{theorem}
\newtheorem{definition}[definition]{Definition}
\aliascntresetthe{definition}
\providecommand{\definitionautorefname}{Definition}
\theoremstyle{remark}
\newtheorem*{remark}{Remark}
\numberwithin{equation}{section}
\begin{document}
\section{Instant theorem and equation previews}
\begin{definition}[Norm]\label{def:norm}
For a real number $x$, its norm is $\lvert x\rvert$.
\end{definition}

\begin{theorem}[Pythagorean identity]\label{thm:identity}
For every real $x$, we have $\sin^2 x+\cos^2 x=1$.
\end{theorem}

\begin{proof}
This follows from the unit circle.
\end{proof}

\begin{align}
 \sin^2 x+\cos^2 x &= 1 \label{eq:identity}\\
 \hat{x} &= \frac{x}{\lVert x\rVert} \tag{N}\label{eq:normal}
\end{align}

Theorem~\ref{thm:identity} gives \eqref{eq:identity}.
The normalization is \eqref{eq:normal}; see also \autoref{def:norm}.

\begin{remark}
Try the Equation ribbon to insert a matrix, an accent, or paired delimiters.
\end{remark}

\section{Practice theorem editing}
\begin{definition}[A positive sequence]\label{def:positive}
A sequence $(a_n)$ is positive if $a_n>0$ for every $n$.

This second paragraph is part of the definition. Press Enter in its body to
add a paragraph using a blank source line.
\end{definition}

\begin{lemma}[A simple consequence]
If $(a_n)$ is positive, then $a_n+a_{n+1}>0$.
\end{lemma}
\begin{proof}
Add the two strict inequalities from Definition~\ref{def:positive}.
\end{proof}

Click a theorem title to add, edit or remove it. The Theorems tab appears only
inside a theorem; Manage theorem environments edits name, printed title, style,
counter sharing and numbering. Home, Paragraph also lists the declared theorem
styles. Try a custom unnumbered remark and inspect its declaration in the preamble.

\section{Prepare appendices}
Appendices use ordinary headings after \verb|\appendix|. In PasTeX, choose
Insert, Environment, Start appendices, then insert a Section heading from Home.
The command changes numbering for the following sections. Sections before it keep their original numbers.

The supporting calculation is in Appendix~\ref{app:calculation}; the extra
assumptions are in Appendix~\ref{app:assumptions}. References, equations and
theorems use the appendix letter automatically.
The command \verb|\autoref| also supplies the name: see \autoref{app:calculation}.

\appendix
\section{Supporting calculation}\label{app:calculation}
\subsection{A finite sum}\label{app:sum}
\begin{equation}\label{eq:appendix-sum}
  \sum_{k=1}^{n} k = \frac{n(n+1)}{2}.
\end{equation}
\begin{theorem}[Induction]
Equation~\eqref{eq:appendix-sum} holds for every positive integer $n$.
\end{theorem}
\begin{proof}
The case $n=1$ is immediate. Add $n+1$ to both sides for the next case.
\end{proof}

\section{Additional assumptions}\label{app:assumptions}
For a book or report, use \verb|\chapter{...}| after \verb|\appendix| instead:
chapters become A, B, and so on, and their sections become A.1, A.2, and so on.
For a document split into files, put \verb|\appendix| before the appendix's
\verb|\input| or \verb|\include| command in the master file.
\end{document}
